% a01_H0rel part 2: the operator layer on the cube and its sub-domains
\subsection{Operator bounds on the cube and its sub-domains}\label{app:a01:ops}

The statements of this subsection are \cite[Thm.~3.3]{K1} and Ba{\l}aban's propagator bounds \cite[Sects.~1--2]{Ba4} for the cube operator $\Mop(\Omega,A)$ of \cref{app:a01:weights}, on every sub-domain $\Omega$ (any union of $j$-blocks of the cube, meeting the Dirichlet planes or not, re-entrant edges allowed, Neumann or Dirichlet cuts on $\dN\Omega$), at every scale $1\le j\le K_f$, uniformly in $\vth$, $J$, $k$, the shape of $\Omega$ and the regular field $A$. At $j=0$ (fixed spacing $5^{-1}$) every clause holds with $L$-dependent constants by finite dimensionality. The constants depend on $(d,L,a,\Wt,c_V)$ and King's constants only.

\begin{lemma}[coercivity]\label{lem:a01:A1}
There are $\gamma=\gamma(d,a,\Wt)>0$ and $\theta_0>0$ such that $\Mop(\Omega,A)\ge\gamma$ for every $\Omega$, $\vth\le1/5$, $V_j\ge0$ and $A$ regular with $\theta_j\le\theta_0$.
\end{lemma}
\begin{proof}
Dropping the bonds between distinct blocks lowers the form (each dropped term is $\ge0$; Dirichlet terms only add), and $P(A)$ is block diagonal (Q8), so $\Mop\ge\bigoplus_B\Mop_B$ with $\Mop_B\defeq-\Delta^N_{A,B}+a_jP_B(A)$ on $\ell^2(B\cap\Xl)$, the killed sites of a face block entering as a Dirichlet term. At $A=0$ write $\phi=c+\psi$ on $B$ with $\psi$ orthogonal to constants (after extending $\phi$ by $0$ to the killed sites, which does not change $P\phi$ by (W3) and only raises the gradient form). The Neumann gap of the unit cube gives $\ip\psi{-\Delta^N\psi}\ge g_P\norm\psi^2$ with $g_P=4\vth^{-2}\sin^2(\pi\vth/2)\ge8$ for $\vth\le\frac12$; by (Q1), (Q4), $|Q\phi|=|c+Q\psi|\ge|c|-\Wt^{1/2}\norm\psi$. Hence $\ip\phi{\Mop_B\phi}\ge g_P\norm\psi^2+a(1-L^{-2})(|c|-\Wt^{1/2}\norm\psi)_+^2\ge\gamma_*(|c|^2+\norm\psi^2)$ with $\gamma_*=\gamma_*(a,\Wt)>0$ by the springs-in-series minimisation $\min_{s\ge0}[\Wt^{-1}s^2+a'(F_0-s)_+^2]=F_0^2/(a'^{-1}+\Wt)$; explicitly $\gamma\ge(1-L^{-2})\min\{a/2,8/(1+2\Wt)\}$. For regular $A$: on the simply connected block, $A=A(x_B)+A''$ with $A(x_B)$ constant, pure gauge, removed by a unitary gauge transformation that preserves both terms, and $|A''|\le\sqrt3K_j$; then $\norm{D_{A''}\phi-D_0\phi}\le2\theta_j\norm\phi$ and $|Q(A'')\phi-Q(0)\phi|\le3\sqrt3\,\Wt^{1/2}\theta_j\norm\phi$ (contours of length $\le3$), so $\ip\phi{\Mop_B\phi}\ge(\gamma_*/2-C\theta_j)\norm\phi^2$, which is the claim for $\theta_j\le\theta_0$.
\end{proof}

\begin{lemma}[Combes--Thomas; unit-lattice kernels]\label{lem:a01:A23}
(a) For every $1$-Lipschitz $g$, $\norm{e^{\delta_0g}\Mop^{-1}e^{-\delta_0g}}_{2\to2}\le2/\gamma$, hence $\norm{\mathbf 1_X\Mop^{-1}\mathbf 1_Y}\le(2/\gamma)e^{-\delta_0\dist(X,Y)}$, with $\delta_0=\delta_0(\gamma,a,\Wt,d)$. (b) $|\Delta^{(j)}(\Omega,A)(u,v)|\le Ce^{-\delta_0|u-v|}$. (c) If $\Omega$ is a union of $(j+1)$-blocks, then $\gamma_0''\le\Delta^{(j)}+aL^{-2}Q_c^*Q_c\le a_j+a$ on the blocks of $\Omega$, with $\gamma_0''=\gamma_0''(d,L,a,\Wt)>0$ for $A=0$ and, at regular $A$ with $\theta_j\le\theta_0'$, half of it; consequently $|C^{(j)}(u,v)|\le Ce^{-\delta_1|u-v|}$.
\end{lemma}
\begin{proof}
(a) Conjugation multiplies the hopping entry $(x,x')$ by $e^{\delta(g(x)-g(x'))}$, $|g(x)-g(x')|\le\vth$; the first-order part is anti-Hermitian and drops from $\Re\ip\phi{\Mop_\delta\phi}$, the remainder has norm $\le3d\delta^2$; the block term changes by $\le(e^{\delta\sqrt3}-1)\Wt\le2\sqrt3\,\Wt\delta$ (its Schur norm, (Q6)); Dirichlet terms are diagonal and commute. So $\Re\Mop_\delta\ge\gamma-3d\delta^2-2\sqrt3\Wt a\delta\ge\gamma/2$ for $\delta\le\delta_0$. (b) $\Delta^{(j)}(u,v)=a_j\delta_{uv}-a_j^2\ip{Q^*e_u}{\Mop^{-1}Q^*e_v}$, $\norm{Q^*e_u}=W_2^{1/2}$ (Q4), and (a). (c) The upper bound is $\Delta^{(j)}\le a_j$, from the minimiser identity $\ip\psi{\Delta^{(j)}\psi}=\min_\Phi[\ip\Phi{H\Phi}+a_j\norm{\psi-Q\Phi}^2]$, $H\defeq-\Delta_A+V_j$ (take $\Phi=0$). Lower bound: dropping the fine bonds between distinct $(j+1)$-blocks $Y$ lowers $H$ and hence $\Delta^{(j)}$ (monotonicity of the minimiser identity), and $Q_c^*Q_c$ is $(j+1)$-block diagonal, so it suffices to bound $F_Y(\psi)\defeq\ip\psi{(\Delta^{(j)}_Y+aL^{-2}Q_c^*Q_c)\psi}$ from below on the $L^3$ coarse sites of one $(j+1)$-block $Y$ (Neumann on its faces, Dirichlet at killed sites). $Q_c^*Q_c$ is the projector onto constants on $Y$; write $\psi=c+\psi_\perp$. Since $Q\mathbf 1=\mathbf 1$ (Q1), $\norm{\psi-Q\Phi}^2\ge\norm{\psi_\perp-(Q\Phi)_\perp}^2\ge(\norm{\psi_\perp}-\norm{Q(\Phi-\av\Phi)})_+^2\ge(\norm{\psi_\perp}-(\Wt C_Ps)^{1/2})_+^2$ when $\ip\Phi{-\Delta^N_Y\Phi}\le s$, by (Q4) and the Neumann Poincar\'e inequality on the live sites of $Y$ (a cube of side $L$, possibly minus a killed face layer, whose Dirichlet bonds only add; constant $C_P(L)$ uniform in $\vth$). Springs in series: $\Delta^{(j)}_Y(\psi)\ge\norm{\psi_\perp}^2/(\Wt C_P+1/a_j)$, and the coarse term is $aL^{-2}|c|^2L^3$; so $F_Y\ge\min\{(\Wt C_P+1/a)^{-1},aL^{-2}\}\norm\psi^2$ at $A=0$. At regular $A$ gauge away the constant part on $Y$ (diameter $\sqrt3L$) at cost $C(L)\theta_j$. The decay of $C^{(j)}$ then follows from the lower bound and (b) by Combes--Thomas for exponentially decaying unit-lattice kernels (the conjugation argument of (a) with the Schur bound $\norm{e^{tg}Xe^{-tg}-X}\le Ct\sum_{z\in\Z^3}|z|e^{-\delta_0|z|/2}$), which is \cite[(5.6)--(5.7)]{Ba4}.
\end{proof}

\begin{lemma}[Kato, and sup bounds with decay]\label{lem:a01:A4}
Let $u=\Mop(\Omega,A)^{-1}f$. (a) $(-\Delta^{N,D}_{0,\Omega}+V_j)|u|\le|f-a_jP(A)u|$ pointwise on $\Omega\cap\Xl$, for every $A$ and $V_j\ge0$ (the free Laplacian with the same cuts and the same Dirichlet terms). (b) For block data $b$ on the blocks of $\Omega$ and the leg $u=a_j\Mop^{-1}Q^*b$,
\[
|u(x)|\le C\sum_{v}e^{-\delta_2\dist(x,B(v))}|b(v)|\le C|b|_\infty,\qquad |(Qu)(v)|\le C\sum_{v'}e^{-\delta_0|v-v'|}|b(v')|,
\]
at every point of every $\Omega$ and every $A$ (given \cref{lem:a01:A1}). (c) The same with $Q^*b$ replaced by any $f$ with block-$L^2$ data $\norm{f\mathbf 1_{B(v)}}_2\le|b(v)|$; in particular $\norm{\Mop^{-1}f}_\infty\le C\norm f_\infty$ and $|\Mop^{-1}f(x)|\le Ce^{-\delta_2\dist(x,\supp f)/2}\norm f_\infty$ for bounded $f$.
\end{lemma}
\begin{proof}
(a) At $u(x)=0$ the left side is $\le0$. For $u(x)\ne0$, $\Re[\bar u(x)(u(x)-U_{xx'}u(x'))]/|u(x)|\ge|u(x)|-|u(x')|$ since $U_{xx'}$ is orthogonal; a face bond contributes $\vth^{-2}|u(x)|$ on both sides; sum over the kept bonds and use $(-\Delta_A+V_j)u=f-a_jPu$. (b), (c) Block averages first: $\norm{\mathbf 1_{B(v)}u}_2\le(2a/\gamma)\sum_{v'}e^{-\delta_0\dist}\norm{\mathbf 1_{B(v')}f}_2$ by \cref{lem:a01:A23}(a), and $|(Qu)(v)|\le\Wt^{1/2}\norm{\mathbf 1_{B(v)}u}_2$ (Q4). Then with $\Mop_1\defeq-\Delta^{N,D}_{0,\Omega}+V_j+a_j$ write $u=\Mop_1^{-1}[f+a_j(1-P)u]$ in the sense of (a): $|u|\le\Mop_1^{-1}|f+a_j(1-P)u|$ by inverse positivity of the $M$-matrix $\Mop_1$, with $\norm{\mathbf 1_B(1-P)u}_2\le\norm{\mathbf 1_Bu}_2$ by (Q3). The kernel of $\Mop_1^{-1}$ is dominated entrywise by that of the Neumann Yukawa operator $(-\Delta^N_{\Omega'}+a_j)^{-1}$ on the union $\Omega'$ of the blocks of $\Omega$ and of the blocks adjacent to its top face blocks (so that the killed sites adjacent to $\Omega\cap\Xl$ are sites of $\Omega'$): $\Mop_1$ is a principal compression of that $M$-matrix plus a non-negative diagonal, and for non-singular $M$-matrices the inverse of a principal compression is dominated entrywise by the compression of the inverse (Schur complement, $(M^{-1})_{SS}=(M_{SS}-M_{SR}M_{RR}^{-1}M_{RS})^{-1}$ with $M_{SR}M_{RR}^{-1}M_{RS}\ge0$ entrywise). \Cref{lem:a01:A4N} gives the block summability $\sup_x\sum_Be^{\delta_2'\dist(x,B)}\norm{K^N_a(x,\cdot)}_{L^2(B)}\le C$ of the Neumann Yukawa kernel on every union of blocks, uniformly in $\vth$; combining the two decays through $\dist(x,B')\le\dist(x,B)+\sqrt3+\dist(B,B')$ gives (b) and (c) with $\delta_2\le\delta_0$. For bounded $f$ only the blocks meeting $\supp f$ carry data.
\end{proof}

\begin{lemma}[block summability of the Neumann Yukawa kernel]\label{lem:a01:A4N}
Let $\Omega'$ be a union of unit blocks of $\vth\Z^3$ (any shape) with Neumann conditions and $K^N_a\defeq(-\Delta^N_{\Omega'}+a_j)^{-1}$, with kernel for the measure $\vth^3$. There are $C_Y',\delta_2'>0$ depending on $a$ only such that, uniformly in $\vth\le1$ and $\Omega'$, $\sup_x\sum_Be^{\delta_2'\dist(x,B)}\norm{K^N_a(x,\cdot)}_{L^2(B)}\le C_Y'$ and $K^N_a(x,y)\le C_Y'e^{-2\delta_2'|x-y|}$ for $|x-y|\ge\frac12$.
\end{lemma}
\begin{proof}
$-\Delta^N_{\Omega'}$ generates the reflecting continuous-time simple random walk on $\Omega'$, so $K^N_a(x,y)=\int_0^\infty e^{-a_jt}p_t(x,y)\,dt$ with $p_t$ symmetric and $\norm{K^N_a(x,\cdot)}_2^2=\int_0^\infty e^{-a_ju}u\,p_u(x,x)\,du$. \emph{On-diagonal bound.} On one unit cube $Q$ the discrete Nash inequality $\norm f_{2,Q}^{2+4/d}\le C_N(E_Q(f)+\norm f_{2,Q}^2)\norm f_{1,Q}^{4/d}$ holds with $C_N$ independent of $\vth$ (the discretisations of one fixed convex domain); it globalises to $\Omega'$ with no geometric hypothesis: with $a_Q\defeq\norm f_{2,Q}^2$, $e_Q\defeq E_Q(f)+a_Q$, $b_Q\defeq\norm f_{1,Q}$ the cube inequality reads $a_Q\le C_N^{d/(d+2)}e_Q^{d/(d+2)}b_Q^{4/(d+2)}$, and summing over the cubes with H\"older's inequality at the exponents $\frac{d+2}d,\frac{d+2}2$, then $\sum b_Q^2\le(\sum b_Q)^2$ and $\sum_QE_Q\le E_{\Omega'}$, gives the Nash inequality on $\ell^2(\Omega')$. By \cite[Thm.~2.1]{CKS87} this is equivalent to $p_u(x,x)\le Cu^{-d/2}$ for $u\le1$; trivially $p_u(x,x)\le\vth^{-d}$ and $p_u\le p_1$ for $u\ge1$. In $d=3$ the integral $\int_0^{\vth^2}\vth^{-3}u\,du+\int_{\vth^2}^1Cu^{-1/2}du+\int_1^\infty e^{-a_ju}Cu\,du$ is bounded, which bounds the near blocks. \emph{Far points}, $D\defeq|x-y|\ge\frac12$: for a $1$-Lipschitz $g$ the conjugated form satisfies $\Re\,e^{\theta g}(-\Delta^N_{\Omega'})e^{-\theta g}\ge-\lambda(\theta)$, $\lambda(\theta)\defeq2d\vth^{-2}(\cosh(\theta\vth)-1)\le d\theta^2\cosh(\theta\vth)$ (bond by bond, each site in $\le2d$ bonds), so $\norm{\mathbf 1_Fe^{t\Delta^N}\mathbf 1_E}\le\exp(-\theta D+d\theta^2t\cosh(\theta\vth))$ for sets at distance $D$, which is $\le e^{-D^2/(5dt)}$ for $t\ge\vth D$ and $\le e^{-D/(8d\vth)}$ for $t\le\vth D$. For $t\le\vth D$ this gives $p_t(x,y)\le\vth^{-d}e^{-D/(8d\vth)}$ and $\int_0^{\vth D}e^{-a_jt}p_t\,dt\le Ce^{-D/(16d)}$. For $\vth D\le t\le1$, Davies' method \cite{Dav93} (the twisted $L^p$ energy inequalities for the conjugated semigroup, with the Nash inequality above and $\lambda_p(\theta)\le C_p\theta^2\cosh(\theta\vth)$, give the $2\to\infty$ bound and duality the $1\to2$ bound) gives $p_t(x,y)\le Ct^{-d/2}e^{-\theta D+C_0dt\theta^2\cosh(\theta\vth)}$, and $\theta=D/(2C_0dt)$ gives $p_t\le Ct^{-d/2}e^{-cD^2/t}$, whose time integral is $\le Ce^{-cD^2/2}$. For $t\ge1$ write $p_t(x,y)=\ip{p_{1/2}(y,\cdot)}{e^{(t-1)\Delta^N}p_{1/2}(x,\cdot)}$, split each $p_{1/2}$ into its part within $D/3$ of its base point and the rest; the main term is $\le Ce^{-D^2/(45dt)}$ by the Davies--Gaffney bound with $\norm{p_{1/2}(x,\cdot)}_2^2=p_1(x,x)\le C$, and the other terms are $\le Ce^{-cD^2}$; then $\int_1^\infty e^{-a_jt}e^{-c_1D^2/t}dt\le(2C/a_j)e^{-\sqrt{2a_jc_1}D}$. Altogether $K^N_a(x,y)\le C_Y'e^{-2\delta_2'D}$ with $2\delta_2'=\min\{\frac1{48d},\frac c2,\sqrt{\frac32ac_1}\}$, and the block sum follows (at most $3^d$ near blocks; far blocks by the pointwise bound and unit volume).
\end{proof}

\begin{lemma}[linear vanishing at the Dirichlet face]\label{lem:a01:DV}
For every $\Omega$, every $A$ (given \cref{lem:a01:A1}), block data $b$ on the blocks of $\Omega$ and every $x\in\Omega\cap\Xl$: $|u(x)|\le C_{\mathrm{DV}}|b|_\infty\min(1,t(x))$ for the leg $u$; the same for $u=\Mop^{-1}f$ with bounded $f$ ($C_{\mathrm{DV}}\norm f_\infty$).
\end{lemma}
\begin{proof}
Fix the face $\mu$ with $t_\mu(x)=t(x)\le1$ (otherwise \cref{lem:a01:A4}). Let $F_0\defeq a_j\Wt(1+C)|b|_\infty$ bound $|a_jQ^*(b-Qu)|$ ((Q6), \cref{lem:a01:A4} for $|Qu|$) and $S_0$ the sup bound of \cref{lem:a01:A4}. Put $\rho(s)\defeq(S_0+F_0)s-F_0s^2/2$ on $[0,1+\vth]$: increasing there, $\rho(0)=0$, $\rho(1)\ge S_0$, $\rho\ge0$, and $-\Delta_s\rho=F_0$ exactly (second difference of a quadratic). Comparison region $R\defeq$ the first-layer slab $\Omega\cap\Xl\cap\{0<t_\mu\le1\}$, with $\Phi(x)\defeq\rho(t_\mu(x))$. On $R$: $(-\Delta^{N,D}_{0,\Omega}+V_j)\Phi\ge F_0\ge$ the right-hand side of Kato (\cref{lem:a01:A4}(a)), because tangential bonds and tangential Neumann faces see $\Phi$ constant, a Neumann face of $\Omega$ at the far side of the first-layer block drops a neighbour with larger $\Phi$ (which raises $-\Delta^N\Phi$), and on face bonds the killed neighbour carries $\Phi=\rho(0)=0=|u|$. The boundary of $R$ inside $\Omega\cap\Xl$ consists of sites at $t_\mu\in\{1,1+\vth\}$, where $|u|\le S_0\le\Phi$. The operator $-\Delta^{N,D}_{0,R}+V_j$ with these data is a non-singular $M$-matrix ($V_j>0$), so the maximum principle gives $|u|\le\Phi\le(S_0+F_0)t_\mu$ on $R$. Other Dirichlet faces met by $R$ carry $u=0\le\Phi$. For bounded $f$ take $F_0\defeq\norm f_\infty+a_j\Wt C\norm f_\infty$.
\end{proof}

No hypothesis on the faces of $\Omega$ is used: on the cube every sub-domain carries the Dirichlet condition of the cube at the Dirichlet planes, and the barrier sees the Neumann cuts only through the zero drop of a function constant along them.

\begin{lemma}[tangentially localised amplitude]\label{lem:a01:DVloc}
There are $\delta_5>0$ and $C$ such that for every $\Omega$, every $A$, every block $v$ of $\Omega$, every face $\mu$ and every $x$ in the first-layer slab of $\mu$, the leg $l(\cdot,v)$ of the unit datum on $v$ satisfies $|l(x,v)|\le C(\vth+t_\mu(x))e^{-\delta_5\dist(x,B(v))}$; consequently $|u(x)|\le C|b|_\infty(\vth+t(x))e^{-\delta_5\dist(x,\supp b)/2}$ there, and the same for bounded sources supported in a union of blocks.
\end{lemma}
\begin{proof}
The barrier of \cref{lem:a01:DV} is multiplied by tangential weights: for each $\nu\ne\mu$ put $f_\nu(s)\defeq\exp(-\delta|\tilde\psi(s)-\tilde\psi(v_\nu)|)$ with $\tilde\psi(s)\defeq s+(2\pi)^{-1}\sin(2\pi(s+\frac12+\frac\vth2))$, whose derivative vanishes exactly on the block midplanes $s\in\Z-\vth/2$, so that $f_\nu$ takes equal values at the two sites adjacent to any block midplane and every Neumann cut of $\Omega$ (which lies on a block midplane) drops a bond of zero difference. The tangential directions are intervals $[0,N]$ with killed ends, on which $u=0\le$ barrier, so no periodisation is needed. Term by term, $-\Delta_\nu f_\nu\ge-c_\delta f_\nu$ with $c_\delta=(2\pi\delta+4\delta^2)e^{3\delta}$ (the kink at $v_\nu$ is concave). With $\Phi\defeq\rho(t_\mu)\prod_{\nu\ne\mu}f_\nu$ and $F_0,S_0$ enlarged by a fixed factor, $(-\Delta^{N,D}+V_j)\Phi\ge\prod f_\nu[F_0-2(c_\delta+8\delta)\sup\rho]\ge\frac12F_0\prod f_\nu$ for $\delta\le\delta_5$ small; the source $|a_jQ^*(\delta_v-Ql)|\le F_0e^{-\delta_2\dist(\cdot,B(v))}$ (\cref{lem:a01:A4}) is dominated by $\frac12F_0\prod f_\nu$ for $\delta\le\delta_2/\sqrt3$, and so are the data at $t_\mu\in\{1,1+\vth\}$. The maximum principle gives $|l|\le\Phi$, and $\prod f_\nu\le e^{2\delta/\pi}e^{-\delta\sum_\nu|x_\nu-v_\nu|}$ with $\dist(x,B(v))\le\sum_\nu|x_\nu-v_\nu|+(t_\mu(v)-1)_++3$. For bounded sources decompose $f=\sum_vf\mathbf 1_{B(v)}$, apply the barrier block by block and sum, halving the rate.
\end{proof}

\paragraph{Interior estimates.} For regular $A$, $r\in[2\vth,1]$ and a ball $B_{2r}(x_0)$ of sites on which $-\Delta_Au=g$ holds as a full-lattice equation, $\sup_{B_r}|D_Au|\le C(r^{-1}\sup_{B_{2r}}|u|+r\sup_{B_{2r}}|g|)$, together with the $\alpha$-H\"older bound of $D_Au$ on $B_r$; this is \cite[(2.31)--(2.33)]{Ba4}, which gauges away the constant part of $A$ on the ball and perturbs from the free lattice Green's function. On the cube the full-lattice equation holds at every live site whose neighbours are live or killed (extend $u$ by $0$ to $\Kill$: a face bond carries no link), so the estimate applies to every ball containing no killed site and no site within one bond of $\dN\Omega$; the open ball $B_{t(x)}(x)$ contains no killed site.

\begin{lemma}[first differences with decay, up to the face]\label{lem:a01:A5}
For every $\Omega$ and regular $A$, for the leg $u$ of data $b$ supported in a set $S$ of blocks, and for $u=\Mop^{-1}f$ with bounded $f$ supported in $S$: $|D_Au(b')|\le C\norm b_\infty\,e^{-\delta_3\dist(b',S)}$ (resp.\ $C\norm f_\infty$) at every bond $b'$ of $\Omega\cap\Xl$ whose endpoints have depth $\ge\frac14$, face bonds included; on the whole cube at every bond. The $\alpha$-H\"older quotients of $D_Au$ on the balls used in the proof obey the same bounds with the same decay.
\end{lemma}
\begin{proof}
Let $x$ be the base of $b'$ and $D\defeq\dist(x,S)$. For $\vth>\frac1{32}$ every clause follows from $|D_Au(b')|\le2\vth^{-1}\max_{x'\sim x}|u(x')|$ and \cref{lem:a01:A4}; let $\vth\le\frac1{32}$. (1) Face bonds: $|D_Au|=|u(x)|/\vth\le C_{\mathrm{DV}}\norm b$ by \cref{lem:a01:DV}, with decay by \cref{lem:a01:DVloc}. (2) $t(x)\le4\vth$: $|D_Au(b')|\le2\vth^{-1}\max_{x'\sim x}|u(x')|\le Ce^{-\delta_5D/2}$ by \cref{lem:a01:DVloc} (amplitude $\vth+t\le6\vth$). (3) $4\vth<t(x)\le\frac18$: $r\defeq t(x)/2\ge2\vth$; $B_{2r}(x)=B_{t(x)}(x)$ has no killed site and, as $\dep(x)\ge\frac14\ge t(x)+\frac18$, no site within one bond of $\dN\Omega$; the source $g\defeq-V_ju-a_jPu+a_jQ^*b$ (resp.\ $f-V_ju-a_jPu$) obeys $|g|\le Ce^{-\delta_5(D-t)/2}$ on the ball, and $\sup_{B_t}|u|\le Ct\,e^{-\delta_5(D-t)/2}$ by \cref{lem:a01:DVloc} (the sites of $B_t(x)$ have $t_\mu\le2t\le\frac14$); the interior estimate gives $|D_Au(b')|\le C(r^{-1}t+r)e^{-\delta_5D/2}\le Ce^{-\delta_5D/2}$ and $[D_Au]_{\alpha;B_r}\le Ct^{-\alpha}e^{-\delta_5D/2}$. (4) $t(x)>\frac18$: $r\defeq\frac1{16}$, $B_{1/8}(x)$ has no killed site and no site within one bond of $\dN\Omega$ (depth $\ge\frac14\ge\frac18+2\vth$); the interior estimate with \cref{lem:a01:A4} gives $|D_Au|,[D_Au]_\alpha\le Ce^{-\delta_2(D-1/8)}$. Take $\delta_3\defeq\delta_5/4$.
\end{proof}

\begin{corollary}[external legs]\label{lem:a01:legs}
With the amplitude profile $\rho_j(x)\defeq\min\{1,\vth+t(x)\}$ and Ba{\l}aban's external-field norm \cite[(1.32)]{Ba3} (sup with the amplitude, first differences, $\alpha_0$-H\"older quotients of the first differences on unit balls), for base points of depth $\ge2$ on $\Omega$ (every base point on the whole cube): (a) $\norm{a_j\Mop^{-1}Q^*b}\le C\norm b_\infty$ for the leg of block data; (b) $\norm{G_{(j)}(\cdot,z)}\le C$ for the slice lines $G_{(j)}=G_j-G_{j-1}$ on unions of $(j+1)$-blocks, $j\ge1$.
\end{corollary}
\begin{proof}
(a) The sup clause with the amplitude $\rho_j^{1/2}$ from \cref{lem:a01:DV,lem:a01:A4} by the geometric mean $\min(P,Q)\le P^{1/2}Q^{1/2}$; the first-difference clause from \cref{lem:a01:A5}; the H\"older clause: for $|x-y|\ge r(x)$ the quotient is $\le2\sup_{B_1(x)}|D_Au|r^{-\alpha_0}$ and $\rho^{1/2}r^{-\alpha_0}\le C$ in each case of the proof of \cref{lem:a01:A5} ($\rho\le2t$, $r=t/2$, $\alpha_0\le\frac12$; or $r=\frac18$; or $\rho\le5\vth$, $r\asymp\vth$); for $|x-y|<r(x)$ the H\"older bounds of that proof give $\rho^{1/2}t^{-\alpha_0}\le C$. (b) The slice identity $G_{(j)}=a_jG_jQ^*C^{(j)}QG_j$ of \cite[(2.17)]{K1}, which is algebra valid for every weight, with \cref{lem:a01:A23,lem:a01:A4}.
\end{proof}
